\documentclass[11pt]{article}
\usepackage[margin=1in]{geometry}
\usepackage{amsmath,amssymb,amsthm,mathtools}
\usepackage{graphicx}
\usepackage[hidelinks]{hyperref}
\usepackage{microtype}

\newtheorem{theorem}{Theorem}
\newtheorem{lemma}{Lemma}
\newtheorem{remark}{Remark}
\newcommand{\Hh}{\mathcal H}
\newcommand{\Kp}{\mathcal K_+}
\newcommand{\B}{\mathcal B}
\newcommand{\ran}{\operatorname{ran}}
\newcommand{\spann}{\operatorname{span}}

\title{Inverse-free Gram block dilations for every $\mathcal C_A$ operator}
\author{Candidate full answer to Question 5.9 of arXiv:2302.13873}
\date{11 August 2026}

\begin{document}
\maketitle

\begin{center}
\fbox{\parbox{0.92\linewidth}{\textbf{Status.}
Candidate full answer, likely valid, pending independent expert review.
The construction covers singular Toeplitz Gram matrices and makes no
closed-range or invertibility assumption.}}
\end{center}

\section{Source question}

Bhat, Ghatak, and Pamula, \emph{Operator moment dilations as block
operators}, arXiv:2302.13873, define $T\in\B(\Hh)$ to be in
$\mathcal C_A$, for a positive invertible $A\in\B(\Hh)$, when the
sequence
\[
 A^{-1/2}T^nA^{-1/2}\qquad(n\geq1)
\]
admits a unitary dilation.  At the end of Section 5 they ask
\cite[Question 5.9]{BGP}:

\begin{quote}
How to write block decompositions for isometric and unitary dilations of
general operators of the $\mathcal C_A$ class?
\end{quote}

\begin{center}
\includegraphics[width=0.96\linewidth]{figures/open_problem_crop.png}
\end{center}

The source gives recurrence formulae for general moment sequences only when
several defect blocks are boundedly invertible.  The point below is that a
finite Gram-space model replaces every such inverse by a quotient and an
orthogonal complement.

\section{Proof intuition}

The moments already specify the inner products that the formal vectors
$V^jh$ must have.  For each $n$, put the first $n+1$ formal copies of $\Hh$
into the Hilbert space obtained from their Toeplitz Gram matrix.  Appending a
zero and prepending a zero are isometries between consecutive Gram spaces;
the latter maps are compatible with the former.  Their inductive-limit shift
is therefore the required minimal isometry.

Orthogonalizing the nested finite Gram spaces gives shells
$\mathcal D_n$.  Since the shift sends the first $n+1$ moments into the first
$n+2$, its matrix on these shells is upper Hessenberg, and every block in
column $n$ is computed inside the finite Gram space of order $n+1$.  This is
the desired inverse-free block algorithm.  Finally, adjoining copies of the
wandering defect $\Kp\ominus V\Kp$ produces a negative identity-shift tail;
one coupling column and the Hessenberg positive corner give the unitary block
matrix.

\section{The block construction}

We use inner products linear in the second variable.  Put
\[
 \zeta(0)=I,\qquad
 \zeta(n)=A^{-1/2}T^nA^{-1/2}\ (n>0),\qquad
 \zeta(-n)=\zeta(n)^*.
\]
For $n\geq0$ define the Toeplitz Gram operator
\[
 G_n=\bigl[\zeta(j-i)\bigr]_{i,j=0}^n\in\B(\Hh^{,n+1}).
\]
If $T\in\mathcal C_A$, then $G_n\geq0$: indeed, for any unitary
dilation $U_0$ of the moments and $x_0,\ldots,x_n\in\Hh$,
\[
 \sum_{i,j=0}^n\langle x_i,\zeta(j-i)x_j\rangle
   =\left\|\sum_{j=0}^nU_0^jx_j\right\|^2.
\]

Let $E_n$ be the completion of $\Hh^{n+1}$ modulo the radical of the
Gram form
\begin{equation}\label{eq:gramform}
 \langle x,y\rangle_n
 =\sum_{i,j=0}^n\langle x_i,\zeta(j-i)y_j\rangle.
\end{equation}
Equivalently, $E_n$ is canonically $\overline{\ran G_n^{1/2}}$ with
the ambient Hilbert-space norm.  The quotient-completion formulation is important: it is
valid even when $G_n$ is singular and $\ran G_n^{1/2}$ is not closed.

On representatives define
\[
 J_n(x_0,\ldots,x_n)=(x_0,\ldots,x_n,0),\qquad
 S_n(x_0,\ldots,x_n)=(0,x_0,\ldots,x_n).
\]
Both descend to isometries $E_n\to E_{n+1}$.  For $J_n$ this follows
from the upper-left principal Gram block, and for $S_n$ from the
lower-right principal block and Toeplitz invariance.  Moreover
\begin{equation}\label{eq:compatibility}
 S_{n+1}J_n=J_{n+1}S_n.
\end{equation}

Let $\Kp=\varinjlim(E_n,J_n)$, with canonical isometries
$\kappa_n:E_n\to\Kp$.  By \eqref{eq:compatibility},
\begin{equation}\label{eq:Vdef}
 V\kappa_nx=\kappa_{n+1}S_nx
\end{equation}
defines an isometry on $\Kp$.  Write $e_jh$ for the vector represented
by $h$ in coordinate $j$.  Then
\begin{equation}\label{eq:moments}
 e_jh=V^je_0h,\qquad e_i^*e_j=\zeta(j-i),\qquad
 e_0^*V^ne_0=\zeta(n).
\end{equation}
In particular, $e_0$ embeds $\Hh$ isometrically and $V$ is an
isometric dilation of the $\mathcal C_A$ moment sequence.

Set
\[
 \mathcal K_n=\kappa_nE_n
   =\overline{\spann}\{e_0\Hh,\ldots,e_n\Hh\},
 \qquad
 \mathcal D_0=\mathcal K_0,
 \quad
 \mathcal D_n=\mathcal K_n\ominus\mathcal K_{n-1}\quad(n\geq1).
\]
Thus
\begin{equation}\label{eq:shellsum}
 \Kp=\bigoplus_{n\geq0}\mathcal D_n.
\end{equation}
Let $Q_i$ be the projection onto $\mathcal D_i$ and put
\begin{equation}\label{eq:blocks}
 V_{ij}=Q_iV|_{\mathcal D_j}:\mathcal D_j\longrightarrow\mathcal D_i.
\end{equation}
Because $V\mathcal K_j\subseteq\mathcal K_{j+1}$,
\begin{equation}\label{eq:hessenbergzero}
 V_{ij}=0\qquad(i>j+1).
\end{equation}
Hence, relative to \eqref{eq:shellsum},
\begin{equation}\label{eq:hessenberg}
 V=\begin{bmatrix}
 V_{00}&V_{01}&V_{02}&V_{03}&\cdots\\
 V_{10}&V_{11}&V_{12}&V_{13}&\cdots\\
 0&V_{21}&V_{22}&V_{23}&\cdots\\
 0&0&V_{32}&V_{33}&\cdots\\
 \vdots&\vdots&\vdots&\ddots&\ddots
 \end{bmatrix}.
\end{equation}

This is also a finite algorithm for the blocks.  To compute column $j$,
construct $E_{j+1}$ from $G_{j+1}$, regard the earlier $E_i$ as the
nested closed ranges of the $J_i$, take their successive orthogonal
differences $\mathcal D_i$, and compress $S_j$ as in
\eqref{eq:blocks}.  Thus column $j$ uses only $G_{j+1}$, orthogonal
complements, and positive square roots.  No inverse, pseudoinverse, or
closed-range hypothesis occurs.

The first column recovers the familiar defect block.  If
$X=\zeta(1)=A^{-1/2}TA^{-1/2}$ and
$D_X=(I-X^*X)^{1/2}$, then the map
\[
 e_1h-e_0Xh\longmapsto D_Xh
\]
extends to a unitary $\mathcal D_1\simeq\overline{\ran D_X}$.  In
these coordinates
\begin{equation}\label{eq:firstblocks}
 V_{00}=X,\qquad V_{10}=D_X.
\end{equation}
The next moment gives an inverse-free characterization of the first
off-diagonal row block:
\begin{equation}\label{eq:nextblock}
 V_{01}D_X=\zeta(2)-X^2.
\end{equation}
Existence and boundedness follow from the Gram construction, and uniqueness
holds because $\ran D_X$ is dense in $\overline{\ran D_X}$.  The later
blocks are obtained by exactly the same finite Gram orthogonalization.

\section{Full theorem and proof}

\begin{theorem}[General $\mathcal C_A$ block dilations]\label{thm:main}
Let $A>0$ be invertible and let $T\in\mathcal C_A$.  The construction
above has the following properties.
\begin{enumerate}
\item The operator $V$ of \eqref{eq:Vdef} is a minimal isometric
dilation of the moment sequence
\[
 \bigl\{A^{-1/2}T^nA^{-1/2}:n\geq1\bigr\}.
\]
Relative to the canonical Gram shells $\mathcal D_n$, its blocks are
given by \eqref{eq:blocks} and have the upper-Hessenberg form
\eqref{eq:hessenberg}.  Every block in column $j$ is determined by the
finite Toeplitz Gram matrix $G_{j+1}$.
\item Put $\mathfrak D=\Kp\ominus V\Kp$ and
\[
 \widehat{\mathcal K}=\Kp\oplus\mathfrak D_{-1}
       \oplus\mathfrak D_{-2}\oplus\cdots,
 \qquad \mathfrak D_{-r}\simeq\mathfrak D.
\]
If $\iota_*:\mathfrak D\hookrightarrow\Kp$ is the inclusion, then
\begin{equation}\label{eq:Uformula}
 \widehat U(x,d_{-1},d_{-2},\ldots)
   =(Vx+\iota_*d_{-1},d_{-2},d_{-3},\ldots)
\end{equation}
is a minimal unitary dilation of the same moments.
\end{enumerate}
Both constructions are valid without any nonsingularity or closed-range
condition on a defect or Gram operator.
\end{theorem}

\begin{proof}
Positivity of the Gram forms, the isometric stage maps, their compatibility,
and the moment identities were proved above.  Equation
\eqref{eq:moments} also gives
\[
 \Kp=\overline{\spann}\{V^ne_0h:n\geq0,\ h\in\Hh\},
\]
so $V$ is minimal.  The shell decomposition and
$V\mathcal K_j\subseteq\mathcal K_{j+1}$ prove the asserted block form.
All objects used for column $j$ lie in $E_{j+1}$, proving finite
determination.

Since $V$ is an isometry,
\[
 \Kp=V\Kp\oplus\mathfrak D.
\]
Consequently the two summands in the first coordinate on the right of
\eqref{eq:Uformula} are orthogonal.  The formula therefore preserves norms.
It is onto as well: given $(y,e_{-1},e_{-2},\ldots)$, decompose uniquely
$y=Vx+\iota_*d$ and take the preimage
$(x,d,e_{-1},e_{-2},\ldots)$.  Thus $\widehat U$ is unitary.  For
$h\in\Hh\subseteq\Kp$ and $n\geq0$, its negative-tail coordinates
remain zero, so $\widehat U^nh=V^nh$.  Equation \eqref{eq:moments}
proves the dilation identity.

It remains only to check two-sided minimality.  The forward iterates of
$\Hh$ span $\Kp$.  Also $\mathfrak D\subseteq\Kp$, and for
$d\in\mathfrak D$,
\[
 \widehat U^{-1}(\iota_*d,0,0,\ldots)=(0,d,0,\ldots).
\]
Approximating $\iota_*d$ by forward orbit vectors shows that the first
negative copy of $\mathfrak D$ lies in the closed two-sided orbit of
$\Hh$.  Repeated application of $\widehat U^{-1}$ gives every further
negative copy.  Hence the closed span of
$\{\widehat U^nh:n\in\mathbb Z,h\in\Hh\}$ is all of
$\widehat{\mathcal K}$.
\end{proof}

For completeness, the unitary matrix can be read directly from
\eqref{eq:Uformula}.  Order the spaces as
\[
 \cdots\oplus\mathfrak D_{-3}\oplus\mathfrak D_{-2}
 \oplus\mathfrak D_{-1}\ \big|\
 \mathcal D_0\oplus\mathcal D_1\oplus\cdots.
\]
There are identity blocks from $\mathfrak D_{-r}$ to
$\mathfrak D_{-(r-1)}$ for $r\geq2$; the column from
$\mathfrak D_{-1}$ into the positive side is
\[
 C_i=Q_i\iota_*:\mathfrak D\to\mathcal D_i\qquad(i\geq0);
\]
and the positive-positive corner is exactly the Hessenberg matrix
\eqref{eq:hessenberg}.  All other negative-to-positive and
positive-to-negative blocks vanish.  This is the requested unitary block
decomposition.

\section{Verification, scope, and novelty}

\paragraph{Degeneracy audit.}
The quotient by the radical makes every $E_n$ well defined when $G_n$ has a
kernel.  Completing the quotient handles nonclosed $\ran G_n^{1/2}$.
Isometric images of the complete stages are closed, so every shell and every
orthogonal projection used above exists.  If a shell or the wandering defect
is zero, the same formula simply omits that summand.

\paragraph{Exact scope.}
The theorem answers the source question in canonical Gram coordinates for
every $T\in\mathcal C_A$.  It does not claim that the blocks can always be
written as bounded rational expressions on fixed copies of $\Hh$; such
expressions are precisely what can fail when the source recurrence asks for
an inverse of a singular defect.  Instead it gives an intrinsic finite-stage
algorithm and a fully explicit action formula.

\paragraph{Bounded novelty check.}
On 11 August 2026 we searched the run's registry and problem indexes for
arXiv:2302.13873 and the core phrases ``$\mathcal C_A$ block dilation'',
``operator-valued Toeplitz moment block Hessenberg'', and the exact wording of
Question 5.9.  We also performed bounded title/citation and web searches for
later papers referring to the question.  No later explicit answer was found.
The positive-kernel/Kolmogorov construction and the minimal unitary extension
of an isometry are classical, and the source itself discusses the abstract
kernel route.  The potentially new observation is the inverse-free finite
Gram-shell extraction that turns this route into the requested general block
model.  Novelty confidence is therefore moderate, while proof-validity
confidence is high.

\paragraph{Human-review recommendation.}
Check (i) the compatibility of $J_n$ and $S_n$ on quotient-completions,
(ii) whether the canonical Gram-shell algorithm matches the intended meaning
of ``write block decompositions'' in Question 5.9, and (iii) the minimality of
the defect-tail unitary.  No computational evidence is used in the proof.

\begin{thebibliography}{9}
\bibitem{BGP}
B. V. Rajarama Bhat, A. Ghatak, and S. K. Pamula,
\emph{Operator moment dilations as block operators},
arXiv:2302.13873 (2023), Question 5.9.
\end{thebibliography}

\end{document}
